\section{PART I --- ORDINARY DIFFERENTIAL EQUATIONS}

\subsection{1. The Geometric View}
\begin{itemize}
\item Vector fields and integral curves; flows and the group property
\item Existence and uniqueness (Picard-Lindel\"of); Peano without uniqueness
\item Continuous dependence on initial conditions and parameters
\item Maximal intervals of existence; conditions for global existence vs.\ blowup
\end{itemize}

\subsection{2. Phase Plane Analysis}
\begin{itemize}
\item Equilibria, stability, basins of attraction
\item Nullclines and phase portraits; qualitative analysis
\item Linearization at equilibria; Jacobian
\item Classification of fixed points: node, spiral, saddle, center
\end{itemize}

\subsection{3. Linear Systems}
\begin{itemize}
\item Matrix exponential $e^{tA}$; computation via Jordan form
\item Fundamental matrix solutions; Wronskian
\item Variation of parameters
\item Floquet theory for periodic systems
\end{itemize}

\subsection{4. Stability Theory}
\begin{itemize}
\item Lyapunov's direct method; Lyapunov functions; LaSalle invariance
\item Stable, unstable, center manifolds (statements + intuition)
\item Hartman-Grobman linearization theorem
\item Input-to-state stability (brief)
\end{itemize}

\subsection{5. Planar Dynamics}
\begin{itemize}
\item Poincar\'e-Bendixson theorem; limit cycles
\item Van der Pol oscillator; relaxation oscillations
\item Index theory (brief); Dulac's criterion (statement)
\item Applications: FitzHugh-Nagumo neuron model; Lotka-Volterra ecology
\end{itemize}

\subsection{6. Bifurcation Theory}
\begin{itemize}
\item Saddle-node, transcritical, pitchfork bifurcations
\item Hopf bifurcation: sub- and supercritical
\item Center manifold reduction; normal forms
\item Codimension-2 bifurcations (brief); applications to bifurcation diagrams
\end{itemize}

\subsection{7. Hamiltonian and Lagrangian Systems}
\begin{itemize}
\item Symplectic structure; Hamiltonian vector fields; Poisson brackets
\item Conservation laws; Noether's theorem
\item Integrable systems (brief); action-angle variables
\item Near-integrable systems; KAM theorem (statement and conceptual meaning only)
\end{itemize}

\subsection{8. Perturbation Methods}
\begin{itemize}
\item Regular perturbation; asymptotic expansions
\item Singular perturbation: matched asymptotics; boundary layers
\item Prandtl boundary layers in fluid mechanics
\item Homogenization in composite materials
\item Method of multiple scales; averaging theory
\item WKB approximation; connections to quantum mechanics
\end{itemize}

\subsection{9. Delay Differential Equations}
\begin{itemize}
\item Constant and state-dependent delays; method of steps
\item Characteristic quasi-polynomial; stability analysis for linear DDEs
\item Hopf bifurcation in DDEs; infinite-dimensional center manifold
\item Applications: population biology; control with delays; neural field equations
\end{itemize}

\subsection{10. Introduction to SDEs}
\begin{itemize}
\item Brownian motion and stochastic processes
\item It\^o integral and It\^o's formula (for SDEs)
\item Basic SDEs: Ornstein-Uhlenbeck, geometric Brownian motion
\item Fokker-Planck equation (connection to Part III)
\item Numerical: Euler-Maruyama; Milstein method
\end{itemize}

\section{PART II --- FUNCTIONAL ANALYTIC FOUNDATIONS}

\subsection{11. Function Spaces for PDEs}
\begin{itemize}
\item $L^p$ spaces; Banach and Hilbert spaces (reprise; now in PDE context)
\item Dense subsets; approximation; compactness in function spaces
\item Weak convergence; weak sequential compactness
\end{itemize}

\subsection{12. Distributions and Weak Derivatives}
\begin{itemize}
\item Test functions; mollifiers and regularization
\item Distributions as linear functionals; Dirac $\delta$ done properly
\item Weak derivatives; operations on distributions
\item Fundamental solutions; examples (Laplacian, heat, wave)
\end{itemize}

\subsection{13. Sobolev Spaces for PDEs}
\begin{itemize}
\item Weak derivatives; $W^{k,p}$ and $H^k$ spaces
\item Sobolev embeddings; trace theorem; Poincar\'e inequality
\item Rellich-Kondrachov compactness; applications to existence theory
\item Fractional Sobolev spaces (brief); boundary regularity
\end{itemize}

\subsection{14. Operator Theory for PDEs}
\begin{itemize}
\item Compact operators; Fredholm alternative in PDE context
\item Spectral theorem for compact self-adjoint operators; eigenfunctions
\item $C_0$-semigroups; abstract Cauchy problem
\item Hille-Yosida theorem; connection to evolution PDEs
\end{itemize}

\section{PART III --- PARTIAL DIFFERENTIAL EQUATIONS}

\subsection{15. Classification and First-Order Equations}
\begin{itemize}
\item Elliptic/parabolic/hyperbolic classification; principal part
\item Method of characteristics for first-order equations; Cauchy problem
\item Eikonal equation; Hamilton-Jacobi (intro)
\item Characteristics for systems; Cauchy-Kovalevskaya (statement)
\end{itemize}

\subsection{16. Elliptic Equations: Weak Formulation and Regularity}
\begin{itemize}
\item Variational formulation of BVPs; energy functional
\item Lax-Milgram theorem; Galerkin method
\item Fredholm alternative for elliptic problems
\item Regularity: interior Schauder, $L^p$ estimates; elliptic bootstrapping
\end{itemize}

\subsection{17. Maximum Principles}
\begin{itemize}
\item Classical maximum principle: elliptic and parabolic versions
\item Comparison principles; Hopf boundary point lemma
\item Alexandroff-Bakelman-Pucci estimate (brief)
\item Applications: uniqueness, a priori bounds, blowup prevention
\end{itemize}

\subsection{18. Spectral Theory for Elliptic Operators}
\begin{itemize}
\item Eigenvalue problems on bounded domains; variational characterization
\item Rayleigh quotient; min-max (Courant-Fischer)
\item Completeness of eigenfunctions; expansion of solutions
\item Weyl's asymptotic law for eigenvalue distribution
\end{itemize}

\subsection{19. Parabolic Equations}
\begin{itemize}
\item Heat kernel; fundamental solution; maximum principle
\item Energy estimates; parabolic regularity
\item Semigroup formulation; analytic semigroups
\item Long-time behavior; connections to probability (Brownian motion)
\end{itemize}

\subsection{20. Hyperbolic Equations}
\begin{itemize}
\item Wave equation: d'Alembert; Kirchhoff; energy conservation
\item Finite speed of propagation; domain of dependence
\item Symmetric hyperbolic systems; energy estimates
\item Conservation laws; weak solutions; Rankine-Hugoniot; entropy conditions
\end{itemize}

\subsection{21. Nonlinear PDEs}
\begin{itemize}
\item Fixed-point methods: Banach contraction, Schauder, Leray-Schauder
\item Variational methods: direct minimization, mountain pass theorem
\item Reaction-diffusion systems; Turing instability; pattern formation
\item Viscosity solutions for Hamilton-Jacobi equations
\end{itemize}

\subsection{22. Inverse Problems for PDEs}
\begin{itemize}
\item Ill-posedness: Hadamard's three conditions; why it arises
\item Tikhonov regularization; regularization parameter selection
\item Parameter identification; coefficient recovery; adjoint method for gradient computation
\item Data assimilation: Kalman filter, ensemble Kalman filter (EnKF), 4D-Var
\item Bayesian inverse problems; posterior; MCMC for PDE-constrained inference
\item Applications: medical imaging, geophysics, materials discovery
\end{itemize}

\subsection{23. Wasserstein Geometry and Mean Field Equations}
\begin{itemize}
\item Optimal transport as a PDE tool; Wasserstein distances
\item Wasserstein gradient flows; JKO scheme; density functional theory
\item McKean-Vlasov SDEs; propagation of chaos
\item Mean field games: MFG system; Nash equilibria; connections to HJB
\item Connection to generative models (diffusion, flow matching)
\end{itemize}

\section{PART IV --- CONTROL THEORY}

\subsection{24. Linear Control Theory}
\begin{itemize}
\item Controllability and observability; Kalman rank conditions
\item Linear quadratic regulator (LQR): cost functional, Riccati equation
\item Pole placement and state feedback; stabilizability
\item Observers; Luenberger observer; separation principle
\end{itemize}

\subsection{25. Optimal Control}
\begin{itemize}
\item Pontryagin maximum principle; Hamiltonian and adjoint equations
\item Hamilton-Jacobi-Bellman equation; dynamic programming
\item Verification theorem; connection to viscosity solutions
\item Minimum time problems; singular arcs (brief)
\end{itemize}

\subsection{26. Nonlinear Control and Feedback Stabilization}
\begin{itemize}
\item Lyapunov-based control design; control Lyapunov functions
\item Passivity and dissipativity; port-Hamiltonian framework
\item Energy-based modeling; power balance; dirac structures
\item Input-output stability; small gain theorem
\end{itemize}

\subsection{27. Connections to Reinforcement Learning}
\begin{itemize}
\item Markov decision processes as discrete optimal control
\item Bellman equation; Q-learning; actor-critic
\item Policy gradient = stochastic optimal control
\item Differential games; multi-agent connections
\end{itemize}

\section{PART V --- NUMERICAL METHODS}

\subsection{28. Numerical ODEs}
\begin{itemize}
\item Runge-Kutta methods; order conditions; Butcher tableaux
\item Stiff equations; A-stability; implicit methods (BDF, Rosenbrock)
\item Symplectic integrators: St\"ormer-Verlet, leapfrog; preservation of structure
\item Adaptive step control; event detection; software (DifferentialEquations.jl)
\end{itemize}

\subsection{29. Classical Numerical PDE Methods}
\begin{itemize}
\item Finite differences: elliptic, parabolic, hyperbolic; von Neumann stability; CFL
\item Finite elements: Galerkin, stiffness matrix, error estimates in Sobolev norms
\item Spectral and pseudospectral methods; exponential convergence for smooth problems
\item Software: FEniCS/Firedrake (FEM), Dedalus (spectral)
\end{itemize}

\subsection{30. Physics-Informed and Data-Driven Methods}
\begin{itemize}
\item PINNs: residual minimization; architecture and training; adaptive collocation
\item Neural operators: DeepONet, Fourier Neural Operator; operator learning framework
\item Failure modes: spectral bias, out-of-distribution, regularity requirements
\item Differentiable solvers; adjoint method for PDE-constrained optimization
\item Hybrid methods: physics-constrained ML; learned preconditioners
\item Uncertainty quantification in data-driven PDE solving
\end{itemize}
